\begin{figure}[htbp]
    \centering
    \resizebox{0.99\linewidth}{!}{%
    \begin{tikzpicture}[
        >=Latex,
        every node/.style={font=\small},
        block/.style={
            draw,
            rounded corners,
            align=center,
            minimum height=11mm,
            text width=25mm,
            inner sep=2pt,
            fill=RoyalBlue!8
        },
        learned/.style={
            block,
            fill=ForestGreen!11
        },
        operation/.style={
            block,
            text width=24mm,
            fill=ForestGreen!11
        },
        graphnode/.style={
            circle,
            draw=ForestGreen!70!black,
            fill=ForestGreen!18,
            minimum size=5mm,
            inner sep=0pt
        },
        targetnode/.style={
            graphnode,
            draw=Orange!80!black,
            fill=Orange!35,
            very thick
        },
        arrow/.style={->, thick},
        graph/.style={ForestGreen!65!black, thick},
        residual/.style={->, thick, dashed, RoyalBlue!75!black}
    ]
        \node[block] (input) at (-6.2,0)
            {current node\\representations \(\mathbf{h}_i\)};
        \node[learned] (coords) at (-3.75,1.25)
            {learned coordinates\\\(\mathbf{s}_i=\mathbf{W}_s\mathbf{h}_i+\mathbf{b}_s\)};
        \node[learned] (features) at (-3.75,-1.25)
            {propagated features\\\(\mathbf{f}_i=\mathbf{W}_f\mathbf{h}_i+\mathbf{b}_f\)};

        \draw[rounded corners, fill=ForestGreen!4, draw=ForestGreen!55!black]
            (-2.05,0.15) rectangle (1.15,2.35);
        \node[font=\small\bfseries] at (-0.45,2.05)
            {learned \(k\)-NN graph};
        \node[targetnode] (gi) at (-0.45,1.12) {\(i\)};
        \node[graphnode] (g1) at (-1.55,1.67) {};
        \node[graphnode] (g2) at (0.55,1.75) {};
        \node[graphnode] (g3) at (-1.35,0.55) {};
        \node[graphnode] (g4) at (0.65,0.55) {};
        \draw[graph, ->] (g1) -- (gi);
        \draw[graph, ->, line width=1.1pt] (g2) -- (gi);
        \draw[graph, ->, line width=1.7pt] (g3) -- (gi);
        \draw[graph, ->, line width=0.7pt] (g4) -- (gi);

        \node[operation] (messages) at (2.85,0)
            {distance-weighted messages\\
             \(\mathbf{m}_{j\rightarrow i}=w_{ij}\mathbf{f}_j\)\\[-1mm]
             {\scriptsize \(w_{ij}=e^{-10d_{ij}^{2}}\)}};
        \node[operation] (aggregate) at (5.75,0)
            {element-wise\\mean and maximum\\
             \([\mathbf{a}^{\mathrm{mean}}_i
             \Vert\mathbf{a}^{\mathrm{max}}_i]\)};
        \node[block, text width=21mm] (output) at (8.35,0)
            {updated node\\representation \(\mathbf{h}'_i\)};
        \node[block] (self) at (2.85,-2.0)
            {linear transform of\\the original \(\mathbf{h}_i\)};

        \draw[arrow] (input) -- (coords);
        \draw[arrow] (input) -- (features);
        \draw[arrow] (coords) -- (-2.05,1.25);
        \draw[arrow] (1.15,1.25) -| (messages.north);
        \draw[arrow] (features) -- (messages);
        \draw[arrow] (messages) -- (aggregate);
        \draw[arrow] (aggregate) -- (output);
        \draw[residual] (input.south) |- (self.west);
        \draw[residual] (self.east) -| (output.south);
    \end{tikzpicture}%
    }
    \caption{Data flow through one \texttt{GravNetConv} operation. Each node
    is projected into a learned coordinate space and connected to its
    \(k\) nearest neighbours within the event. Neighbour features are weighted
    by learned-space distance, aggregated by their element-wise mean and
    maximum, and combined with a transformed copy of the original node.
    Original diagram based on the GravNet operation introduced in
    Ref.~\cite{qasim2019-gravnet}.}
    \label{fig:gravnet_convolution}
\end{figure}
